A subsequent full-device reference search separates physical opportunity from fixed-other-net capacity feasibility. After pruning only the target-exclusive chain, the reverse query terminates at retained-tree anchors and charges their physical prefix once. It admits any physical crossing count and excludes retained IPIN/SINK nodes as anchors. Thus it is an exact local zero-RC graph reference, not a claim about the historical native bounding box, queued seeds, or global routing/STA optimum. The capacity adapter passes 2,000 independent forward-Floyd comparisons on 1,000 random directed graphs (751 reachable and 1,249 unreachable cases). Physical RRG edges, all 1,028,568 RR capacities, ownership and preserved branches are independently checked.

For net 6590 pin 3, the relaxed reference is 0.856978\,ns, while the fixed-other-net optimum equals the current 1.042804\,ns. For net 5479 pin 36, the relaxed reference is 0.663513\,ns and the capacity-feasible reference is 1.309816\,ns, versus the current 3.265075\,ns. The latter establishes a candidate improvement of 1.955259\,ns within the stated local model. Fixed-capacity query core times are 0.008217 and 0.064732\,s; these exclude graph loading, preparation and validation.

The shorter net 5479 witness is then serialized into an isolated route copy. Two unchanged-VPR analysis-only loads compare the original and repaired routes using identical packed/place/RRG inputs, switch-ID verification, native FULL checks and full STA. Both pass. Independent physical-delay comparison covers 23,905 sink connections: only net 5479 pin 36 changes; all other net traces are identical. Whole-design CPD improves only from 6.20812 to 6.20393\,ns (0.067492\%), as the worst path moves from bram\_select[7] to bram\_rdata\_ext[15], to bram\_select[6] to bram\_rdata\_ext[60]. This supports feasible local repair and critical-path migration, not a 1.955259\,ns CPD gain. The repaired CPD remains worse than the matched count-two A0/R1 result of 6.01810\,ns.

Planar wirelength remains 228,420, but occupied SB pairs increase from 118 to 119 and CB crossing output sources from 955 to 956. Analysis-only process times are 8.586962 and 9.784481\,s; they are load/check/STA costs, not full route-process timings or a formal speed comparison. There are two new native analysis calls and zero new route-search calls. A reporting-wrapper path error occurred after both native processes completed; a separate recovery uses the exact graph from their manifests, with no native reruns. Original failures and recovery evidence are retained.

The following registered continuation selects targets from the accepted state's STA, rather than retaining the original worst-path targets. It preserves a transactional best route and native legality/STA acceptance. Worst-path versus near-critical-set selection, fixed versus refreshed targets, and full-device versus native-domain search are distinct factors. Matched A0/R1/Tf must receive the same post-route repair opportunity before any method-quality claim; this retrospective single-circuit repair is not a promoted online algorithm.
